% Prerequisites: foundational infrastructure for the formalization.
% Structure follows template.tex: §0.1 Spectral Sequences, §0.2 Stable Homotopy Theory, §0.3 Synthetic Spectra.

\section{Prerequisites}\label{sec:prereq}

%%%%%%%%%%%%%%%%%%%%%%%%%%%%%%%%%%%%%%%%%%%%%%%%%%%%%%%%%%%%%%%%%%
\subsection{Spectral Sequences}\label{sec:prereq-ss}
%%%%%%%%%%%%%%%%%%%%%%%%%%%%%%%%%%%%%%%%%%%%%%%%%%%%%%%%%%%%%%%%%%

\subsubsection{Spectral Sequences}\label{sec:prereq-ss-def}

\begin{definition}\label{prereq:def:ss-data}
    \lean{KIP.SpectralSequence.SSData}\leanok
    Let $\mathcal{C}$ be an abelian category. The \emph{nested subspace data} (or \emph{SS data}) for a spectral sequence at a single grading index consists of:
    \begin{enumerate}
        \item An ambient object $V \in \mathcal{C}$.
        \item A decreasing family of subobjects $\{Z_r\}_{r \in \mathbb{N} \cup \{\infty\}}$ of $V$, called \emph{$r$-cycles}:
        $$V = Z_0 \supseteq Z_1 \supseteq Z_2 \supseteq \cdots \supseteq Z_\infty.$$
        \item An increasing family of subobjects $\{B_r\}_{r \in \mathbb{N} \cup \{\infty\}}$ of $V$, called \emph{$r$-boundaries}:
        $$B_0 \subseteq B_1 \subseteq B_2 \subseteq \cdots \subseteq B_\infty.$$
        \item The containment $B_r \subseteq Z_r$ for all $r$.
        \item $Z_\infty$ is the greatest lower bound of the finite cycles: for any subobject $X \subseteq V$, if $X \subseteq Z_i$ for all $i \in \mathbb{N}$, then $X \subseteq Z_\infty$. Together with the monotonicity of $Z$, this gives $Z_\infty = \bigcap_{i \in \mathbb{N}} Z_i$.
        \item $B_\infty$ is the least upper bound of the finite boundaries: for any subobject $X \subseteq V$, if $B_i \subseteq X$ for all $i \in \mathbb{N}$, then $B_\infty \subseteq X$. Together with the monotonicity of $B$, this gives $B_\infty = \bigcup_{i \in \mathbb{N}} B_i$.
    \end{enumerate}
    Together these form the nested chain
    $$V \supseteq Z_0 \supseteq Z_1 \supseteq \cdots \supseteq Z_\infty \supseteq B_\infty \supseteq \cdots \supseteq B_1 \supseteq B_0.$$
    The \emph{$r$-th page} at this index is defined as $E_r = Z_r / B_r$. The \emph{$E_\infty$-page} is $E_\infty = Z_\infty / B_\infty$.
\end{definition}

\begin{definition}\label{prereq:def:spectral-sequence}
    \uses{prereq:def:ss-data}
    \lean{KIP.SpectralSequence.SpectralSequence}\leanok
    A \emph{spectral sequence} in an abelian category $\mathcal{C}$ with index type $\iota$ (typically $\mathbb{Z}^n$) consists of:
    \begin{enumerate}
        \item A starting page index $r_0 \in \mathbb{Z}$.
        \item For each $r \ge r_0$ and index $\mathbf{k} \in \iota$, an object $E_r^{\mathbf{k}} \in \mathcal{C}$ (the \emph{$E_r$-page} at index $\mathbf{k}$).
        \item A \emph{differential degree function} $\mathbf{d}: \mathbb{Z} \to \iota$, and differentials
        $$d_r^{\mathbf{k}}: E_r^{\mathbf{k}} \to E_r^{\mathbf{k} + \mathbf{d}(r)}$$
        satisfying $d_r \circ d_r = 0$ componentwise.
        \item For each index $\mathbf{k}$, nested subspace data $(V^{\mathbf{k}}, Z_r^{\mathbf{k}}, B_r^{\mathbf{k}})$ such that $E_r^{\mathbf{k}} = Z_r^{\mathbf{k}} / B_r^{\mathbf{k}}$ for $r \ge r_0$.
        \item An isomorphism $H(E_r, d_r) \cong E_{r+1}$ for each $r$. By the nesting $B_r \subseteq B_{r+1} \subseteq Z_{r+1} \subseteq Z_r$, the differential $d_r$ has kernel $Z_{r+1}/B_r$ and image $B_{r+1}/B_r$, so
        $$H(E_r, d_r) = \frac{Z_{r+1}/B_r}{B_{r+1}/B_r} \cong Z_{r+1}/B_{r+1} = E_{r+1}$$
        by the third isomorphism theorem.
    \end{enumerate}
\end{definition}

\begin{remark}\label{prereq:rem:decreasing-filtration}
    All filtrations considered in this project are decreasing.
\end{remark}

\begin{definition}\label{prereq:def:einfty-data}
    \uses{prereq:def:ss-data}
    \lean{KIP.SpectralSequence.SSData.eInfty}\leanok
    The \emph{$E_\infty$-page} of a spectral sequence $E$ at index $\mathbf{k}$ is given by the nested subspace data: $E_\infty^{\mathbf{k}} = Z_\infty^{\mathbf{k}} / B_\infty^{\mathbf{k}}$. For bounded or degenerate spectral sequences, $E_\infty = E_r$ for all sufficiently large $r$.
\end{definition}

\begin{definition}\label{prereq:def:ss-morphism}
    \uses{prereq:def:spectral-sequence}
    \lean{KIP.SpectralSequence.SpectralSequenceMorphism}\leanok
    A \emph{morphism of spectral sequences} $f: E \to E'$ consists of a family of maps $\varphi_{\mathbf{k}}: V^{\mathbf{k}} \to V'^{\mathbf{k}}$ on the underlying objects at each index, preserving the cycle and boundary subobjects:
    $$\varphi_{\mathbf{k}}(Z_r^{\mathbf{k}}) \subseteq Z'^{\mathbf{k}}_r, \qquad \varphi_{\mathbf{k}}(B_r^{\mathbf{k}}) \subseteq B'^{\mathbf{k}}_r$$
    for all $r$, and commuting with differentials on the induced page maps. A morphism induces maps on all pages $E_r \to E'_r$ and on the $E_\infty$-page.
\end{definition}

\subsubsection{Convergence of Spectral Sequences}\label{sec:prereq-convergence}

\begin{definition}\label{prereq:def:convergence}
    \uses{prereq:def:spectral-sequence,prereq:def:einfty-data,prereq:def:filtration,prereq:def:associated-graded}
    \lean{KIP.SpectralSequence.Convergence}\leanok
    Let $E$ be a spectral sequence with $n$-grading and let $A$ be an $(n-1)$-graded object equipped with a decreasing filtration (Definition~\ref{prereq:def:filtration}). We say $E$ \emph{converges} to $A$, written $E_{r_0} \Longrightarrow A$, if there exists an isomorphism of $n$-graded objects
    $$E_\infty \cong \mathrm{gr}\, A,$$
    where the $n$-grading on $\mathrm{gr}\, A$ is obtained from the $(n-1)$-grading of $A$ together with the filtration index, possibly composed with a fixed reindexing map.
\end{definition}

\begin{remark}\label{prereq:rem:weak-convergence}
    This is convergence in the weak sense: we only require an isomorphism between $E_\infty$ and $\mathrm{gr}\, A$. Conditional convergence and strong convergence (in the sense of Boardman~\cite{Boardman}) are not considered in this project.
\end{remark}

\begin{definition}\label{prereq:def:converging-ss-morphism}
    \uses{prereq:def:convergence,prereq:def:ss-morphism,prereq:def:filtration,prereq:def:filtered-morphism}
    \lean{KIP.SpectralSequence.ConvergenceMorphism}\leanok
    Given a reindexing rule, the \emph{category of converging spectral sequences} has objects $(E, A, F, \varphi)$ where $E$ is a spectral sequence, $A$ is a filtered graded object, and $\varphi: E_\infty \xrightarrow{\sim} \mathrm{gr}\, A$ is the convergence isomorphism. A morphism $(E_1, A_1) \to (E_2, A_2)$ consists of:
    \begin{enumerate}
        \item A morphism of spectral sequences $E_1 \to E_2$ (Definition~\ref{prereq:def:ss-morphism}), inducing maps on $E_\infty$-pages.
        \item A filtration-preserving map $A_1 \to A_2$.
        \item Compatibility: the induced $E_\infty$-map and the associated graded map commute with the convergence isomorphisms.
    \end{enumerate}
\end{definition}

\begin{definition}\label{prereq:def:stem}
    \uses{prereq:def:convergence}
    \lean{KIP.SpectralSequence.Convergence.stemDegree}\leanok
    In the setting of Definition~\ref{prereq:def:convergence}, the component of the $n$-grading on $E$ that corresponds to the original $(n-1)$-grading of $A$ is called the \emph{stem} (or \emph{topological degree}).
\end{definition}

\begin{definition}\label{prereq:def:filtration-degree}
    \uses{prereq:def:convergence}
    \lean{KIP.SpectralSequence.Convergence.filtrationDegree}\leanok
    The component of the $n$-grading on $E$ that corresponds to the filtration on $A$ is called the \emph{filtration degree}.
\end{definition}

\begin{definition}\label{prereq:def:detect}
    \uses{prereq:def:convergence}
    \lean{KIP.SpectralSequence.Detects}\leanok
    Suppose $E$ converges to $A$. Let $x \in A$ be an element with filtration at least $r$, i.e., $x \in F^r A$. Let $\bar{x}$ denote the image of $x$ in $\mathrm{gr}^r A = F^r A / F^{r+1} A$, and let $y \in E_\infty$ be the element corresponding to $\bar{x}$ under the convergence isomorphism. We say $y$ \emph{detects} $x$.
\end{definition}

\begin{proposition}\label{prereq:prop:detect-zero}
    \uses{prereq:def:detect}
    \lean{KIP.SpectralSequence.detect_zero}\leanok
    An element $x \in F^r A$ is detected by $0 \in E_\infty$ if and only if $x \in F^{r+1} A$.
\end{proposition}
\begin{proof}\leanok Immediate from the definition of detection and the convergence isomorphism.\end{proof}

\begin{proposition}\label{prereq:prop:detect-difference}
    \uses{prereq:def:detect}
    \lean{KIP.SpectralSequence.detect_difference}\leanok
    Two elements $x, x' \in F^r A$ are both detected by the same $y \in E_\infty$ if and only if $x - x' \in F^{r+1} A$.
\end{proposition}
\begin{proof}\leanok Follows from Proposition~\ref{prereq:prop:detect-zero} applied to $x - x'$.\end{proof}

\subsubsection{Filtered Complex}\label{sec:prereq-filtered-complex}

\begin{definition}\label{prereq:def:filtration}
    \lean{KIP.SpectralSequence.Filtration}\leanok
    A \emph{filtered $n$-graded abelian group} is an $n$-graded abelian group $A$ equipped with a decreasing filtration: for each index $\mathbf{k} \in \mathbb{Z}^n$, a family of subgroups
    $$\cdots \supseteq F^{s-1} A^{\mathbf{k}} \supseteq F^s A^{\mathbf{k}} \supseteq F^{s+1} A^{\mathbf{k}} \supseteq \cdots$$
\end{definition}

\begin{definition}\label{prereq:def:filtered-morphism}
    \uses{prereq:def:filtration}
    \lean{KIP.SpectralSequence.FilteredMorphism}\leanok
    A \emph{morphism of filtered graded abelian groups} is a graded homomorphism $f: A \to B$ that preserves the filtration: $f(F^s A^{\mathbf{k}}) \subseteq F^s B^{\mathbf{k}}$ for all $s$ and $\mathbf{k}$. Filtered graded abelian groups and their morphisms form a category.
\end{definition}

\begin{definition}\label{prereq:def:filtered-complex}
    \uses{prereq:def:filtration}
    \lean{KIP.SpectralSequence.FilteredComplex}\leanok
    A \emph{filtered chain complex} is a chain complex $(C_*, d)$ in the category of filtered graded abelian groups: the differential $d: C_n \to C_{n-1}$ preserves the filtration, i.e., $d(F^s C_n) \subseteq F^s C_{n-1}$ for all $s$ and $n$.
\end{definition}

\begin{proposition}\label{prereq:prop:homology-filtration}
    \uses{prereq:def:filtered-complex}
    \lean{KIP.SpectralSequence.FilteredComplex.homologyFiltration}\leanok
    A filtered chain complex $(C_*, d, F^s)$ induces a natural filtration on its homology:
    $$F^s H_n(C_*) = \mathrm{im}(H_n(F^s C_*) \to H_n(C_*)) = \ker(H_n(C_*) \to H_n(C_* / F^s C_*)).$$
\end{proposition}

\begin{definition}\label{prereq:def:associated-graded}
    \uses{prereq:def:filtration}
    \lean{KIP.SpectralSequence.Filtration.associatedGraded}\leanok
    Given a filtered graded abelian group $(A, F)$, the \emph{associated graded} is
    $$\mathrm{gr}^s A^{\mathbf{k}} = F^s A^{\mathbf{k}} / F^{s+1} A^{\mathbf{k}}.$$
\end{definition}

\begin{definition}\label{prereq:def:associated-graded-complex}
    \uses{prereq:def:filtered-complex,prereq:def:associated-graded}
    \lean{KIP.SpectralSequence.FilteredComplex.assocGraded}\leanok
    Given a filtered chain complex $(C_*, d, F^s)$, the \emph{associated graded complex} is the bigraded object $\mathrm{gr}^s C_n$ equipped with the induced differential $\bar{d}: \mathrm{gr}^s C_n \to \mathrm{gr}^s C_{n-1}$ (well-defined since $d$ preserves the filtration).
\end{definition}

\begin{definition}\label{prereq:def:exhaustive}
    \uses{prereq:def:filtered-complex}
    \lean{KIP.SpectralSequence.Filtration.IsExhaustive}\leanok
    A filtered chain complex $(C_*, d, F^s)$ is \emph{exhaustive} if $\bigcup_s F^s C_n = C_n$ for all $n$.
\end{definition}

\begin{definition}\label{prereq:def:hausdorff}
    \uses{prereq:def:filtered-complex}
    \lean{KIP.SpectralSequence.Filtration.IsHausdorff}\leanok
    A filtered chain complex $(C_*, d, F^s)$ is \emph{Hausdorff} (or \emph{separated}) if $\bigcap_s F^s C_n = 0$ for all $n$.
\end{definition}

\begin{definition}\label{prereq:def:bounded-filtration}
    \uses{prereq:def:exhaustive,prereq:def:hausdorff}
    \lean{KIP.SpectralSequence.FilteredComplex.IsBounded}\leanok
    A filtered chain complex $(C_*, d, F^s)$ is \emph{bounded} if for each $n$, there exist $a \le b$ with $F^a C_n = C_n$ and $F^{b+1} C_n = 0$. A bounded filtration is both exhaustive and Hausdorff. We assume all filtered complexes in this project are exhaustive and Hausdorff.
\end{definition}

\begin{definition}\label{prereq:def:filtered-complex-ssdata}
    \uses{prereq:def:ss-data,prereq:def:filtered-complex,prereq:def:associated-graded-complex}
    \lean{KIP.SpectralSequence.FilteredComplex.toSSData}\leanok
    Given a filtered chain complex $(C_*, d, F^s)$, we construct an SSData structure. The spectral sequence starts at $E_0$ with $(n+1)$-grading.
    \begin{enumerate}
        \item Define $V^{s,t} = \mathrm{gr}^s C_{s+t}$, the associated graded.
        \item Define $Z_r^{s,t}$ as the image in $\mathrm{gr}^s C_{s+t}$ of those $x \in F^s C_{s+t}$ with $dx \in F^{s+r} C_{s+t-1}$.
        \item Define $B_r^{s,t}$ as the image in $\mathrm{gr}^s C_{s+t}$ of those elements of the form $dy$ with $y \in F^{s-r} C_{s+t+1}$.
    \end{enumerate}
    The differential $d_r: E_r^{s,t} \to E_r^{s+r,t-r+1}$ is induced by $d$: for $[x] \in Z_r^{s,t}$, define $d_r[x] = [dx] \in E_r^{s+r,t-r+1}$.
\end{definition}

\begin{proposition}\label{prereq:prop:filtered-complex-ssdata-wd}
    \uses{prereq:def:filtered-complex-ssdata}
    \lean{KIP.SpectralSequence.FilteredComplex.pageDifferential_comp}\leanok
    The differential $d_r$ in Definition~\ref{prereq:def:filtered-complex-ssdata} is well-defined and satisfies $d_r \circ d_r = 0$.
\end{proposition}
\begin{proof}\leanok This is verified by a direct computation in the filtered complex, showing that the composition of two successive page differentials factors through zero.\end{proof}

\begin{theorem}\label{prereq:thm:filtered-complex-ss}
    \uses{prereq:def:spectral-sequence,prereq:def:filtered-complex-ssdata,prereq:prop:filtered-complex-ssdata-wd}
    \lean{KIP.SpectralSequence.FilteredComplex.toSpectralSequence}\leanok
    The nested subspace data and differentials from Definition~\ref{prereq:def:filtered-complex-ssdata} satisfy the axioms of a spectral sequence. Specifically:
    \begin{enumerate}
        \item $B_r \subseteq B_{r+1} \subseteq Z_{r+1} \subseteq Z_r$ for all $r$.
        \item $E_r^{s,t} = Z_r^{s,t} / B_r^{s,t}$.
        \item $H(E_r, d_r) \cong E_{r+1}$.
    \end{enumerate}
    The spectral sequence starts at $E_0^{s,t} = F^s C_{s+t} / F^{s+1} C_{s+t}$ with differentials $d_r: E_r^{s,t} \to E_r^{s+r,t-r+1}$.
\end{theorem}

\begin{theorem}\label{prereq:thm:filtered-complex-convergence}
    \uses{prereq:thm:filtered-complex-ss,prereq:def:convergence,prereq:prop:homology-filtration,prereq:def:exhaustive,prereq:def:hausdorff}
    \lean{KIP.SpectralSequence.FilteredComplex.weakConvergence}\leanok
    If the filtration on a filtered chain complex is exhaustive and Hausdorff, then the associated spectral sequence converges to the homology:
    $$E_r^{s,t} \Longrightarrow H_{s+t}(C_*).$$
    The convergence isomorphism is $E_\infty^{s,t} \cong \mathrm{gr}^s H_{s+t}(C_*)$ with respect to the induced filtration from Proposition~\ref{prereq:prop:homology-filtration}.
\end{theorem}

\begin{definition}\label{prereq:def:filtered-complex-morphism}
    \uses{prereq:def:filtered-complex}
    A \emph{morphism of filtered chain complexes} $f: (C_*, F) \to (D_*, G)$ is a chain map $f: C_* \to D_*$ preserving the filtration: $f(F^s C_n) \subseteq G^s D_n$ for all $s, n$.
\end{definition}

\begin{proposition}\label{prereq:prop:filtered-complex-functoriality}
    \uses{prereq:thm:filtered-complex-ss,prereq:def:filtered-complex-morphism,prereq:def:ss-morphism}
    A morphism of filtered chain complexes induces a morphism of the associated spectral sequences. This construction is functorial.
\end{proposition}

\subsubsection{Crossing of Differentials}\label{sec:prereq-crossing}

\begin{definition}\label{prereq:def:essential-differential}
    \uses{prereq:def:spectral-sequence}
    \lean{KIP.SpectralSequence.IsEssentialAt}\leanok
    Let $E$ be a spectral sequence. A differential $d_r(x) = y$ is called \emph{essential} if $y \ne 0$ on the $E_r$-page. Equivalently, $d_r(x) = y$ is essential if $x$ is not a boundary and $y$ is not zero in $E_r$. A differential is \emph{inessential} if $y = 0$ on the $E_r$-page.
\end{definition}

\begin{definition}\label{prereq:def:differential-data}
    \uses{prereq:def:spectral-sequence,prereq:def:essential-differential}
    \lean{KIP.SpectralSequence.DifferentialDatum}\leanok
    A \emph{differential datum} in a spectral sequence $E$ is a triple $(r, x, y)$ where $r \ge r_0$, $x \in E_r^{\mathbf{k}}$, and $y \in E_r^{\mathbf{k}+\mathbf{d}(r)}$ with $d_r(x) = y$. We say the differential datum is \emph{essential} if $y \ne 0$, and denote it $d_r(x) = y$.
\end{definition}

\begin{definition}\label{prereq:def:crossing}
    \uses{prereq:def:convergence,prereq:def:essential-differential,prereq:def:differential-data}
    \lean{KIP.SpectralSequence.HasCrossingAt}\leanok
    Let $E$ be a spectral sequence converging to $A$. Let $x \in E$ be an element of filtration $s$, and suppose $d_n(x) = y$ is a differential (so $y$ has filtration $s + n$). Assume the spectral sequence starts at page $\ge 0$ and that $d_n$ increases filtration by $n$.
    \begin{enumerate}
        \item We say that $d_n(x) = y$ \emph{has a crossing that hits filtration $p$} for some $p \le s+n$, if there exists an element $x' \in E$ of filtration $s + a$ with $a > 0$ and an essential differential $d_m(x') = y'$ for $0 \ne y'$ of filtration $s + a + m$, such that
        $$p = \mathrm{Fil}(y') = s + a + m \le \mathrm{Fil}(y) = s + n.$$
        \item We say that $d_n(x) = y$ \emph{has no crossing that hits the range $\mathrm{Fil} \ge p$} if there does not exist $x'$ of filtration $s + a$ with $a > 0$ and an essential differential $d_m(x') = y'$ for $0 \ne y'$ of filtration $s + a + m$ such that
        $$p \le \mathrm{Fil}(y') = s + a + m \le \mathrm{Fil}(y) = s + n.$$
        \item We say that $d_n(x) = y$ \emph{has no crossing} if it has no crossing that hits the range $\mathrm{Fil} \ge \mathrm{Fil}(x) + 1$.
    \end{enumerate}
\end{definition}

\begin{proposition}\label{prereq:prop:no-crossing-iff}
    \uses{prereq:def:crossing}
    \lean{KIP.SpectralSequence.ess_no_crossing_char}\leanok
    An $f$-extension from $x$ to $y$ has no crossing in the range $\mathrm{AF} \ge p$ if and only if for all $[x] \in \{x\}$ such that $\mathrm{AF}(f[x]) \ge p$ we have $f[x] \in \{y\}$. In particular:
    \begin{enumerate}
        \item An $f$-extension from $x$ to $y$ has no crossing if and only if for all $[x] \in \{x\}$ we have $f[x] \in \{y\}$.
        \item An $f$-extension $d_n^f(x) = 0$ has no crossing if and only if for all $[x] \in \{x\}$, $\mathrm{AF}(f[x]) > \mathrm{AF}(x) + n$.
        \item If $\mathrm{AF}(f) = n$, then all $d_n^f$-differentials have no crossing.
    \end{enumerate}
\end{proposition}

\subsubsection{Extension Spectral Sequence}\label{sec:prereq-ess}

\begin{definition}\label{prereq:def:ess}
    \uses{prereq:def:convergence,prereq:def:adams-filtration}
    \lean{KIP.SpectralSequence.ExtensionSS}\leanok
    For a map $f: X \to Y$ between spectra, the \emph{$f$-extension spectral sequence} ($f$-ESS) is the spectral sequence derived from filtering the chain complex
    $$0 \to \pi_* X \xrightarrow{\pi_* f} \pi_* Y \to 0$$
    by the Adams filtration. Its $E_0$-page is isomorphic to $E_\infty^{s,t}(X) \oplus E_\infty^{s,t}(Y)$, and it converges to $\ker(\pi_* f) \oplus \mathrm{coker}(\pi_* f)$. Differentials have the form
    $$d_n^f: \mathrm{ESS}(f)_n^{s,t} \to \mathrm{ESS}(f)_n^{s+n,t+n}.$$
\end{definition}

\begin{definition}\label{prereq:def:f-extension}
    \uses{prereq:def:ess}
    \lean{KIP.SpectralSequence.HasFExtension}\leanok
    We say there is an \emph{$f$-extension} from $x \in E_\infty^{s,t}(X)$ to $y \in E_\infty^{s+n,t+n}(Y)$ if $d_n^f(x) = y$ in the $f$-ESS. The extension is \emph{essential} if $y$ is nontrivial on the $E_n$-page; otherwise it is \emph{inessential}.
\end{definition}

\begin{notation}\label{prereq:not:representatives}
    \uses{prereq:def:detect}
    \lean{KIP.SpectralSequence.DetectionSet}\leanok
    For $x \in E_\infty^{s,t}(X)$, let $\{x\} \subseteq \pi_{t-s}X$ denote the set of elements detected by $x$:
    $$\{x\} = \{ \alpha \in F^s \pi_{t-s}X \mid \bar{\alpha} \mapsto x \text{ under } E_\infty^{s,t} \cong \mathrm{gr}^s \pi_{t-s}X \}.$$
    For $[x] \in \{x\}$, we call $[x]$ a \emph{representative} of $x$. Note that $\{x\}$ is a coset of $F^{s+1}\pi_{t-s}X$ in $F^s\pi_{t-s}X$.
\end{notation}

\begin{proposition}\label{prereq:prop:f-ext-characterization}
    \uses{prereq:def:f-extension,prereq:not:representatives}
    \lean{KIP.SpectralSequence.ess_diff_char}\leanok
    An $f$-extension $d_n^f(x) = y$ holds if and only if there exists $[x] \in \{x\}$ such that $f[x] \in \{y\}$.
\end{proposition}

\begin{corollary}\label{prereq:cor:af-vanishing}
    \uses{prereq:def:ess,prereq:def:adams-filtration-maps}
    \lean{KIP.SpectralSequence.af_vanishing}\leanok
    If $\mathrm{AF}(f) = k$, then $d_i^f = 0$ for $i < k$.
\end{corollary}

\begin{proposition}\label{prereq:prop:ess-exactness}
    \uses{prereq:prop:f-ext-characterization,prereq:cor:composition-zero,prereq:not:representatives}
    \lean{KIP.SpectralSequence.ess_exactness}\leanok
    If $X \xrightarrow{f} Y \xrightarrow{g} Z$ induces an exact sequence on homotopy groups, then all permanent $d^g$-cycles in $E_\infty(Y)$ are boundaries in the $f$-ESS.
\end{proposition}

\subsubsection{Commutativity of Extension Differentials}\label{sec:prereq-commutativity}

\begin{theorem}\label{prereq:thm:four-spectra}
    \uses{prereq:prop:no-crossing-iff}
    \lean{KIP.SpectralSequence.fourSpectra_exact}\leanok
    Consider a homotopy commutative diagram of spectra
    $$\xymatrix{
    X \ar[r]^{f}\ar[d]_{p} & Y\ar[d]^{q}\\
    Z \ar[r]_{g} & W
    }$$
    Suppose $m, n, l \ge 0$, $0 < k \le m + l - n$, and
    \begin{enumerate}
        \item $d_n^f(x) = y$,
        \item $d_m^p(x) = z$,
        \item the differential in (1) or (2) has no crossing,
        \item $d_l^g(z) = w$, and has no crossing that hits $\mathrm{AF} \ge s + n + k$,
        \item $d_{k-1}^q y = 0$, and has no crossing.
    \end{enumerate}
    Then $d_{m+l-n}^q(y) = w$.
\end{theorem}

\begin{corollary}\label{prereq:cor:four-spectra}
    \uses{prereq:thm:four-spectra}
    \lean{KIP.SpectralSequence.fourSpectra_cor}\leanok
    In the setting of Theorem~\ref{prereq:thm:four-spectra}, if condition~(4) is strengthened to ``$d_l^g(z) = w$ has no crossing'' (and condition~(5) is dropped), then $d_{m+l-n}^q(y) = w$.
\end{corollary}

\begin{corollary}\label{prereq:cor:triangle-map}
    \uses{prereq:cor:four-spectra}
    \lean{KIP.SpectralSequence.essCommutativity_triangle}\leanok
    Consider a homotopy commutative diagram $X \xrightarrow{f} Y \xrightarrow{q} Z$ with $p = q \circ f$. If $d_n^f(x) = y$, $d_m^p(x) = z$, and one of the two extensions has no crossing, then $d_{m-n}^q(y) = z$.
\end{corollary}

\begin{corollary}\label{prereq:cor:composition-ext}
    \uses{prereq:cor:four-spectra}
    \lean{KIP.SpectralSequence.ess_composition}\leanok
    For $q = g \circ p$, if $d_m^p(x) = z$ and $d_l^g(z) = w$ has no crossing, then $d_{m+l}^q(x) = w$.
\end{corollary}

\begin{corollary}\label{prereq:cor:composition-zero}
    \uses{prereq:cor:four-spectra,prereq:prop:f-ext-characterization}
    \lean{KIP.SpectralSequence.ess_composition_zero}\leanok
    If $g \circ f = 0$ and $d_n^f(x) = y$, then $y$ is a permanent cycle in the $g$-ESS: $d_m^g(y) = 0$ for all $m \ge 0$.
\end{corollary}

\begin{corollary}\label{prereq:cor:ess-naturality}
    \uses{prereq:def:ess,prereq:cor:four-spectra}
    \lean{KIP.SpectralSequence.essCommutativity_induces_map}\leanok
    Consider a homotopy commutative diagram
    $$\xymatrix{
    X \ar[r]^{f}\ar[d]_{p} & Y\ar[d]^{q}\\
    Z \ar[r]_{g} & W
    }$$
    If $\mathrm{ESS}(p)_0 = \mathrm{ESS}(p)_r$ and $\mathrm{ESS}(q)_0 = \mathrm{ESS}(q)_r$ for some $r \ge 0$, then $(d_r^p, d_r^q)$ induces a map from the $f$-ESS to the $g$-ESS.
\end{corollary}

\begin{corollary}\label{prereq:cor:stationary-page}
    \uses{prereq:thm:four-spectra}
    \lean{KIP.SpectralSequence.essCommutativity_eInfty_compat}\leanok
    In the setting of Theorem~\ref{prereq:thm:four-spectra}, if the $g$-ESS and the $q$-ESS have stationary pages (i.e., $E_0 = E_r$ for all $r$), then the differentials commute with the induced maps on $E_\infty$-pages.
\end{corollary}

\begin{corollary}\label{prereq:cor:vanishing-target}
    \uses{prereq:thm:four-spectra}
    \lean{KIP.SpectralSequence.essCommutativity_null}\leanok
    In the setting of Theorem~\ref{prereq:thm:four-spectra}, if $W = 0$, then $d_n^f(x) = y$ and $d_m^p(x) = z$ implies $d_{m+l-n}^q(y) = 0$ for all $l$, i.e., $y$ is a permanent cycle in the $q$-ESS.
\end{corollary}


%%%%%%%%%%%%%%%%%%%%%%%%%%%%%%%%%%%%%%%%%%%%%%%%%%%%%%%%%%%%%%%%%%
\subsection{Stable Homotopy Theory}\label{sec:prereq-stable}
%%%%%%%%%%%%%%%%%%%%%%%%%%%%%%%%%%%%%%%%%%%%%%%%%%%%%%%%%%%%%%%%%%

\subsubsection{Stable Homotopy Theory}\label{sec:prereq-stable-homotopy}

\begin{definition}\label{prereq:def:triangulated-category}
    \lean{CategoryTheory.Pretriangulated}\leanok
    A \emph{triangulated category} is an additive category $\mathcal{T}$ equipped with:
    \begin{enumerate}
        \item An additive autoequivalence $\Sigma: \mathcal{T} \to \mathcal{T}$, called the \emph{shift} (or \emph{suspension}) functor.
        \item A class of \emph{distinguished triangles} $X \to Y \to Z \to \Sigma X$, satisfying the axioms (TR1)--(TR4):
        \begin{itemize}
            \item[\textup{(TR1)}] Every morphism $f: X \to Y$ extends to a distinguished triangle $X \xrightarrow{f} Y \to Z \to \Sigma X$. The triangle $X \xrightarrow{\mathrm{id}} X \to 0 \to \Sigma X$ is distinguished.
            \item[\textup{(TR2)}] A triangle $X \to Y \to Z \to \Sigma X$ is distinguished if and only if $Y \to Z \to \Sigma X \to \Sigma Y$ is distinguished.
            \item[\textup{(TR3)}] Given a morphism of the first two terms of two distinguished triangles, there exists a (not necessarily unique) morphism of the third terms making the diagram commute.
            \item[\textup{(TR4)}] The octahedral axiom.
        \end{itemize}
    \end{enumerate}
\end{definition}

\begin{definition}\label{prereq:def:cofiber-sequence}
    \uses{prereq:def:triangulated-category}
    \lean{KIP.StableHomotopy.HoCofiberSequence}\leanok
    In a triangulated category, a distinguished triangle $X \xrightarrow{f} Y \to C_f \to \Sigma X$ is called a \emph{cofiber sequence}, and $C_f$ is the \emph{cofiber} (or \emph{cone}) of $f$. By (TR1), every morphism admits a cofiber sequence.
\end{definition}

\begin{definition}\label{prereq:def:functorial-cofiber}
    \uses{prereq:def:cofiber-sequence}
    \lean{KIP.StableHomotopy.HasFunctorialCofiber}\leanok
    A \emph{functorial cofiber} on a triangulated category $\mathcal{T}$ is an assignment that to each morphism $f: X \to Y$ in $\mathcal{T}$ associates a distinguished triangle $X \xrightarrow{f} Y \to Cf \to \Sigma X$ that is functorial: given a commutative square
    $$\xymatrix{
    X \ar[r]^{f}\ar[d]_{\alpha} & Y\ar[d]^{\beta}\\
    X' \ar[r]_{f'} & Y'
    }$$
    there is an induced map $C\alpha\beta: Cf \to Cf'$ making the diagram of distinguished triangles commute.
\end{definition}

\begin{proposition}\label{prereq:prop:les-hom}
    \uses{prereq:def:cofiber-sequence}
    \lean{CategoryTheory.Pretriangulated.Triangle.coyoneda_exact₂}\leanok
    For a distinguished triangle $X \to Y \to Z \to \Sigma X$ and any object $W$, the long exact sequence on hom-sets
    $$\cdots \to [W, X] \to [W, Y] \to [W, Z] \to [W, \Sigma X] \to \cdots$$
    holds (covariant), and dually for the contravariant case.
\end{proposition}

\begin{definition}\label{prereq:def:closed-symmetric-monoidal}
    \lean{CategoryTheory.MonoidalClosed}\leanok
    A \emph{closed symmetric monoidal category} is a category $\mathcal{C}$ equipped with:
    \begin{enumerate}
        \item A unit object $S$.
        \item A bifunctor $(X, Y) \mapsto X \wedge Y$ (the \emph{smash product} or \emph{tensor product}) from $\mathcal{C} \times \mathcal{C}$ to $\mathcal{C}$, which is associative and commutative up to coherent natural isomorphism, with $S \wedge X \cong X$ up to coherent natural isomorphism.
        \item Function objects $F(X, Y)$, functorial contravariantly in $X$ and covariantly in $Y$, satisfying the tensor-hom adjunction
        $$[X, F(Y, Z)] \cong [X \wedge Y, Z]$$
        naturally in all three variables.
    \end{enumerate}
\end{definition}

\begin{definition}\label{prereq:def:closed-symmetric-tensor-triangulated}
    \uses{prereq:def:closed-symmetric-monoidal,prereq:def:triangulated-category}
    \lean{KIP.StableHomotopy.ClosedSymmetricTensorTriangulated}\leanok
    A \emph{closed symmetric tensor triangulated category} is a triangulated category $\mathcal{C}$ equipped with a closed symmetric monoidal structure (Definition~\ref{prereq:def:closed-symmetric-monoidal}) that is \emph{compatible} with the triangulation:
    \begin{enumerate}
        \item The smash product preserves suspensions: there is a natural equivalence $e_{X,Y}: \Sigma X \wedge Y \xrightarrow{\sim} \Sigma(X \wedge Y)$, compatible with the unit and associativity isomorphisms.
        \item The smash product is exact: if $X \xrightarrow{f} Y \xrightarrow{g} Z \xrightarrow{h} \Sigma X$ is an exact triangle, then for any object $W$,
        $$X \wedge W \xrightarrow{f \wedge 1} Y \wedge W \xrightarrow{g \wedge 1} Z \wedge W \xrightarrow{h \wedge 1} \Sigma(X \wedge W)$$
        is exact.
        \item The functor $F(X, Y)$ is exact in $Y$; it is exact in $X$ up to sign.
        \item The smash product interacts with suspension in a graded-commutative manner: the twist map $T: S^r \wedge S^s \to S^s \wedge S^r$ and the equivalence $S^r \wedge S^s \simeq S^{r+s}$ satisfy $T = (-1)^{rs}$.
    \end{enumerate}
    This definition follows Hovey--Palmieri--Strickland~\cite{HPS97}, Definition~A.2.1.
\end{definition}

\begin{definition}\label{prereq:def:tensor-triangulated-functorial-cofiber}
    \uses{prereq:def:closed-symmetric-tensor-triangulated,prereq:def:functorial-cofiber}
    \lean{KIP.StableHomotopy.TensorTriangulatedCatWithFunctorialCofiber}\leanok
    A \emph{closed symmetric tensor triangulated category with functorial cofiber} is a closed symmetric tensor triangulated category $\mathcal{T}$ (Definition~\ref{prereq:def:closed-symmetric-tensor-triangulated}) equipped with:
    \begin{enumerate}
        \item A functorial cofiber (Definition~\ref{prereq:def:functorial-cofiber}).
        \item A \emph{tensor-cofiber exchange}: for morphisms $f: X \to Y$ and $g: A \to B$, a natural isomorphism
        $$C(f \wedge g) \cong Cf \wedge Cg$$
        compatible with the triangulated structure.
    \end{enumerate}
\end{definition}

\begin{definition}\label{prereq:def:stable-homotopy-category}
    \uses{prereq:def:tensor-triangulated-functorial-cofiber}
    \lean{KIP.StableHomotopy.StableHomotopyCategory}\leanok
    The \emph{stable homotopy category} $\mathcal{S}$ is a closed symmetric tensor triangulated category with functorial cofiber (Definition~\ref{prereq:def:tensor-triangulated-functorial-cofiber}). Its objects are called \emph{spectra}. Homotopy pullback squares and homotopy pushout squares in $\mathcal{S}$ coincide.
\end{definition}

\begin{remark}
    In this project, $\mathcal{S}$ refers to the homotopy category of spectra (in the sense of stable homotopy theory). It can be constructed as the homotopy category of a stable model category, but we take it as given and work axiomatically with its triangulated and closed symmetric monoidal structure.
\end{remark}

\begin{definition}\label{prereq:def:sphere-spectrum}
    \uses{prereq:def:stable-homotopy-category}
    \lean{KIP.StableHomotopy.SphereSpectrum}\leanok
    The \emph{sphere spectrum} $\mathbb{S} \in \mathcal{S}$ is the unit object for the smash product.
\end{definition}

\begin{definition}\label{prereq:def:sphere-Sn}
    \uses{prereq:def:sphere-spectrum}
    \lean{KIP.StableHomotopy.Sphere}\leanok
    For an integer $n \in \mathbb{Z}$, define
    $$S^n = \Sigma^n \mathbb{S},$$
    the $n$-fold suspension of the sphere spectrum.
\end{definition}

\begin{definition}\label{prereq:def:homotopy-groups}
    \uses{prereq:def:sphere-Sn}
    \lean{KIP.StableHomotopy.HomotopyGroup}\leanok
    For a spectrum $X \in \mathcal{S}$, the \emph{$n$-th homotopy group} of $X$ is
    $$\pi_n X = [S^n, X] = \Hom_{\mathcal{S}}(S^n, X).$$
\end{definition}

\begin{proposition}\label{prereq:prop:homotopy-graded-abelian}
    \uses{prereq:def:homotopy-groups}
    \lean{KIP.StableHomotopy.HomotopyGroups}\leanok
    For any spectrum $X$, the collection $\pi_* X = \{\pi_n X\}_{n \in \mathbb{Z}}$ is a graded abelian group.
\end{proposition}

\begin{definition}\label{prereq:def:smash-product}
    \uses{prereq:def:stable-homotopy-category,prereq:def:closed-symmetric-tensor-triangulated,prereq:def:sphere-Sn}
    \lean{KIP.StableHomotopy.smashSpheres_iso}\leanok
    The \emph{smash product} $\wedge: \mathcal{S} \times \mathcal{S} \to \mathcal{S}$ is the symmetric monoidal product from the closed symmetric monoidal structure of $\mathcal{S}$ (Definition~\ref{prereq:def:closed-symmetric-tensor-triangulated}), with unit $\mathbb{S}$. It satisfies $S^m \wedge S^n \simeq S^{m+n}$ and induces a pairing on homotopy groups:
    $$\pi_m X \otimes \pi_n Y \to \pi_{m+n}(X \wedge Y).$$
\end{definition}

\begin{definition}\label{prereq:def:mapping-spectra}
    \uses{prereq:def:stable-homotopy-category,prereq:def:closed-symmetric-tensor-triangulated,prereq:def:smash-product}
    \lean{KIP.StableHomotopy.MappingSpectrum}\leanok
    For spectra $X, Y \in \mathcal{S}$, the \emph{mapping spectrum} $F(X, Y)$ is the internal hom (function object) from the closed symmetric monoidal structure (Definition~\ref{prereq:def:closed-symmetric-tensor-triangulated}), adjoint to the smash product:
    $$[X \wedge Y, Z] \cong [X, F(Y, Z)].$$
    In particular, $\pi_n F(X, Y) \cong [\Sigma^n X, Y]$.
\end{definition}

\begin{proposition}\label{prereq:prop:les-homotopy}
    \uses{prereq:def:cofiber-sequence,prereq:def:homotopy-groups}
    \lean{KIP.StableHomotopy.les_homotopy_exact_f}\leanok
    Let $X \xrightarrow{f} Y \to C_f \to \Sigma X$ be a cofiber sequence. Then there is a long exact sequence of homotopy groups:
    $$\cdots \to \pi_{n+1}(C_f) \to \pi_n(X) \xrightarrow{f_*} \pi_n(Y) \to \pi_n(C_f) \to \pi_{n-1}(X) \to \cdots$$
    This follows from Proposition~\ref{prereq:prop:les-hom} and the representability $\pi_n(-) = [S^n, -]$.
\end{proposition}
\begin{proof}\leanok Follows from the long exact Hom sequence applied to representable functors.\end{proof}

\subsubsection{Cohomology}\label{sec:prereq-cohomology}

\begin{definition}\label{prereq:def:eilenberg-maclane}
    \uses{prereq:def:stable-homotopy-category}
    \lean{KIP.StableHomotopy.EilenbergMacLane}\leanok
    The \emph{Eilenberg--MacLane spectrum} $\HF$ is the spectrum representing mod~2 ordinary cohomology. It is a commutative ring spectrum in $\mathcal{S}$.
\end{definition}

\begin{axiom}\label{prereq:ax:eilenberg-maclane-homotopy}
    \uses{prereq:def:eilenberg-maclane,prereq:def:homotopy-groups}
    \lean{KIP.StableHomotopy.HF2_pi0}\leanok
    The homotopy groups of the Eilenberg--MacLane spectrum are
    $$\pi_n(\HF) = \begin{cases} \bF_2 & \text{if } n = 0, \\ 0 & \text{if } n \neq 0.\end{cases}$$
\end{axiom}

\begin{definition}\label{prereq:def:mod2-cohomology}
    \uses{prereq:def:eilenberg-maclane,prereq:def:mapping-spectra}
    \lean{KIP.StableHomotopy.Mod2CohomologyTotal}\leanok
    The \emph{mod~2 cohomology} of a spectrum $X$ is
    $$H^n(X; \bF_2) = [\Sigma^n X, \HF].$$
\end{definition}

\begin{definition}\label{prereq:def:mod2-homology}
    \uses{prereq:def:eilenberg-maclane,prereq:def:smash-product,prereq:def:homotopy-groups}
    \lean{KIP.StableHomotopy.Mod2HomologyTotal}\leanok
    The \emph{mod~2 homology} of a spectrum $X$ is
    $$H_n(X; \bF_2) = \pi_n(\HF \wedge X).$$
\end{definition}

\subsubsection{Adams Spectral Sequence}\label{sec:prereq-adams}

\begin{axiom}\label{prereq:ax:adams-ss}
    \uses{prereq:def:spectral-sequence,prereq:def:stable-homotopy-category,prereq:def:converging-ss-morphism}
    \lean{KIP.StableHomotopy.AdamsSS}\leanok
    The \emph{$\HF$-Adams spectral sequence} is a functor from $\mathcal{S}$ to the category of bigraded converging spectral sequences. For each spectrum $X$, it produces a spectral sequence $\{E_r^{s,t}(X), d_r\}$ with $r \ge 2$ and differentials
    $$d_r: E_r^{s,t} \to E_r^{s+r, t+r-1}.$$
    A map $f: X \to Y$ induces maps $f_*: E_r^{s,t}(X) \to E_r^{s,t}(Y)$ compatible with differentials.
\end{axiom}

\begin{axiom}\label{prereq:ax:adams-convergence}
    \uses{prereq:ax:adams-ss,prereq:def:convergence}
    \lean{KIP.StableHomotopy.adamsConvergence}\leanok
    For a finite spectrum $X$, the Adams spectral sequence converges:
    $$E_2^{s,t}(X) \Longrightarrow \pi_{t-s} X.$$
    More precisely, $E_\infty^{s,t}(X) \cong F^s \pi_{t-s} X / F^{s+1} \pi_{t-s} X$.
\end{axiom}

\begin{axiom}\label{prereq:ax:adams-e2-order2}
    \uses{prereq:ax:adams-ss}
    \lean{KIP.StableHomotopy.adams_2torsion}\leanok
    All elements in the $E_2$-page of the mod~2 Adams spectral sequence have order~2, i.e., $2x = 0$ for all $x \in E_2^{s,t}(X)$. This holds because $E_2^{s,t}(X)$ is an $\bF_2$-vector space.
\end{axiom}

\begin{definition}\label{prereq:def:adams-filtration}
    \uses{prereq:ax:adams-ss,prereq:def:filtration,prereq:def:homotopy-groups}
    \lean{KIP.StableHomotopy.InAdamsFiltration}\leanok
    The Adams spectral sequence induces a natural decreasing filtration on $\pi_* X$, called the \emph{Adams filtration}. An element $\alpha \in \pi_n X$ has \emph{Adams filtration} $\ge s$ if it lies in $F^s \pi_n X$.
\end{definition}

\begin{definition}\label{prereq:def:adams-filtration-maps}
    \uses{prereq:def:adams-filtration,prereq:def:mod2-homology}
    \lean{KIP.StableHomotopy.AF}\leanok
    The Adams filtration extends to maps between spectra. For a map $f: X \to Y$, the \emph{Adams filtration} $\AF(f)$ is the largest $k$ such that $f$ factors as a composition of $k$ maps, each inducing the zero map on $\HF$-homology.
\end{definition}

\begin{proposition}\label{prereq:prop:adams-filtration-characterization}
    \uses{prereq:def:adams-filtration-maps}
    \lean{KIP.StableHomotopy.adamsFiltration_functorial}\leanok
    A map $f: X \to Y$ with $\AF(f) = k$ maps $F^s \pi_* X$ into $F^{s+k} \pi_* Y$ for all $s$.
\end{proposition}


%%%%%%%%%%%%%%%%%%%%%%%%%%%%%%%%%%%%%%%%%%%%%%%%%%%%%%%%%%%%%%%%%%
\subsection{$\HF$-Synthetic Spectra}\label{sec:prereq-synthetic}
%%%%%%%%%%%%%%%%%%%%%%%%%%%%%%%%%%%%%%%%%%%%%%%%%%%%%%%%%%%%%%%%%%

\subsubsection{Synthetic Spectra}\label{sec:prereq-syn-basic}

\begin{definition}\label{prereq:def:synthetic-category}
    \uses{prereq:def:stable-homotopy-category,prereq:def:functorial-cofiber}
    \lean{KIP.Synthetic.SyntheticCategory}\leanok
    The category $h\mathrm{Syn}$ of \emph{$\HF$-synthetic spectra} is a triangulated category with functorial cofiber (it arises as the homotopy category of a stable model category $\mathrm{Syn}_{\HF}$).
\end{definition}

\begin{definition}\label{prereq:def:bigraded-suspension}
    \uses{prereq:def:synthetic-category}
    \lean{KIP.Synthetic.SyntheticCategory.biShift}\leanok
    The category $h\mathrm{Syn}$ admits a \emph{bigraded suspension} $\Sigma^{m,n}$ for each $(m,n) \in \mathbb{Z}^2$, which is an autoequivalence of $h\mathrm{Syn}$ satisfying
    $$\Sigma^{m_1,n_1}\Sigma^{m_2,n_2} \cong \Sigma^{m_1+m_2, n_1+n_2}.$$
    The functor $\Sigma^{1,0}$ is identified with the triangulated suspension functor.
\end{definition}

\begin{definition}\label{prereq:def:lambda}
    \uses{prereq:def:bigraded-suspension}
    \lean{KIP.Synthetic.SyntheticCategory.lam}\leanok
    There exists a natural transformation
    $$\lambda: \Sigma^{0,-1} \to \mathrm{Id}$$
    in $h\mathrm{Syn}$. For a synthetic spectrum $X$, we define $X/\lambda^n$ as the cofiber of $\lambda^n: \Sigma^{0,-n}X \to X$.
\end{definition}

\begin{remark}
    The hom-sets in $h\mathrm{Syn}$ are enriched over $\mathbb{Z}[\lambda]$-modules: $\Hom(X,Y)$ is the bigraded abelian group $[\Sigma^{*,*}X, Y]$, with $\lambda$ acting by the natural transformation of degree $(0,-1)$.
\end{remark}

\subsubsection{Synthetic Spheres}\label{sec:prereq-syn-spheres}

\begin{definition}\label{prereq:def:bigraded-spheres}
    \uses{prereq:def:bigraded-suspension,prereq:def:synthetic-category}
    \lean{KIP.Synthetic.Smn}\leanok
    The \emph{bigraded spheres} are the synthetic spectra $S^{m,n}$ defined by
    $$S^{m,n} = \Sigma^{m,n} S^{0,0}$$
    where $S^{0,0}$ is the sphere object of $h\mathrm{Syn}$.
\end{definition}

\begin{definition}\label{prereq:def:synthetic-homotopy-groups}
    \uses{prereq:def:bigraded-spheres}
    \lean{KIP.Synthetic.BiHom}\leanok
    The \emph{synthetic homotopy groups} of a synthetic spectrum $X$ are
    $$\pi_{m,n}(X) = [S^{m,n}, X].$$
    These form a bigraded abelian group with a natural $\mathbb{Z}[\lambda]$-module structure.
\end{definition}

\begin{proposition}\label{prereq:prop:synthetic-homotopy-shift}
    \uses{prereq:def:synthetic-homotopy-groups,prereq:def:bigraded-suspension}
    \lean{KIP.Synthetic.susp_invariance}\leanok
    $$\pi_{m,n}(X) = \pi_{m+k,n+l}(\Sigma^{k,l}X).$$
\end{proposition}

\subsubsection{Synthetic Adams Spectral Sequence}\label{sec:prereq-syn-adams}

\begin{axiom}\label{prereq:ax:syn-adams-ss}
    \uses{prereq:def:spectral-sequence,prereq:def:synthetic-category,prereq:def:converging-ss-morphism}
    \lean{KIP.Synthetic.SynAdamsSS}\leanok
    The \emph{synthetic Adams spectral sequence} is a functor from $\mathcal{S}$ to the category of 3-graded converging spectral sequences. For a spectrum $X$, it produces a spectral sequence $\SynSS_r^{s,t,w}(X)$ with differentials
    $$d_r: \SynSS_r^{s,t,w} \to \SynSS_r^{s+r, t+r-1, w}.$$
    It converges to $\pi_{*,*}(\nu X)$.
\end{axiom}

\begin{remark}
    The synthetic Adams spectral sequence is a $\mathbb{Z}[\lambda]$-module spectral sequence, with $\lambda$ in tridegree $(0,0,-1)$.
\end{remark}

\subsubsection{The $\nu$ Functor}\label{sec:prereq-nu}

\begin{definition}\label{prereq:def:nu-functor}
    \uses{prereq:def:stable-homotopy-category,prereq:def:synthetic-category}
    \lean{KIP.Synthetic.nu}\leanok
    There exists a functor
    $$\nu: h\mathcal{S} \to h\mathrm{Syn}$$
    from the homotopy category of spectra to the homotopy category of synthetic spectra. The functor $\nu$ is lax symmetric monoidal and preserves filtered colimits.
\end{definition}

\begin{remark}
    \uses{prereq:def:nu-functor}
    The functor $\nu$ does not commute with suspension: instead, $\nu(\Sigma X) \cong \Sigma^{1,1} \nu X$. The natural comparison map $\Sigma(\nu X) \to \nu(\Sigma X)$ is $\lambda$.
\end{remark}

\begin{axiom}\label{prereq:ax:nu-cofiber}
    \uses{prereq:def:nu-functor,prereq:def:cofiber-sequence,prereq:def:mod2-homology}
    \lean{KIP.Synthetic.nu_cofiber_iff}\leanok
    Suppose $X \xrightarrow{f} Y \xrightarrow{g} Z$ is a cofiber sequence of spectra. Then $\nu X \xrightarrow{\nu f} \nu Y \xrightarrow{\nu g} \nu Z$ is a cofiber sequence of synthetic spectra if and only if
    $$0 \to {\HF}_*X \xrightarrow{{\HF}_*f} {\HF}_*Y \xrightarrow{{\HF}_*g} {\HF}_*Z \to 0$$
    is a short exact sequence.
\end{axiom}

\subsubsection{Synthetic Rigidity}\label{sec:prereq-syn-rigidity}

\begin{axiom}\label{prereq:ax:rigid}
    \uses{prereq:def:lambda,prereq:ax:syn-adams-ss}
    \lean{KIP.Synthetic.rigidity_degeneration}\leanok
    The synthetic Adams spectral sequence for $\nu X$ has $E_2$-page
    $$\SynSS_2^{*,*,*}(\nu X) \cong E_2^{*,*}(X) \otimes \bF_2[\lambda],$$
    where an element in $E_2^{s,t}(X)$ has tridegree $(s,t,t)$ and $\lambda$ has tridegree $(0,0,-1)$. Given a classical Adams differential $d_r^{\mathrm{cl}}(x) = y$, the corresponding synthetic differential is $d_r(x) = \lambda^{r-1}y$, which is $\lambda$-linear, and all synthetic Adams differentials arise in this way.
\end{axiom}

\begin{axiom}\label{prereq:ax:lambda-bockstein}
    \uses{prereq:def:lambda,prereq:ax:rigid}
    \lean{KIP.Synthetic.lambda_bockstein_iso}\leanok
    The synthetic Adams spectral sequence for $\nu X$ is isomorphic to the $\lambda$-Bockstein spectral sequence:
    \begin{enumerate}
        \item $E_2^{s,t}(X) \cong \pi_{t-s,t}(\nu X / \lambda)$.
        \item If there is a classical Adams differential $d_r x = y$ for $x \in E_2^{s,t}(X) \cong \pi_{t-s,t}(\nu X / \lambda)$, then $x$ admits a lift to $\pi_{t-s,t}(\nu X / \lambda^{r-1})$ whose image under the Bockstein $\nu X / \lambda^{r-1} \to \Sigma^{1,-r+1} \nu X / \lambda$ equals $d_r(x)$.
    \end{enumerate}
\end{axiom}

\begin{proposition}\label{prereq:prop:syn-einfty-nuX}
    \uses{prereq:ax:rigid}
    \lean{KIP.Synthetic.einfty_nuX}\leanok
    The $E_\infty$-page of the synthetic Adams spectral sequence for $\nu X$ is
    $$E_\infty^{s,t,w}(\nu X) \cong \begin{cases}
        Z_\infty^{s,t}(X) / B_{1+t-w}^{s,t}(X) & \text{if } t \ge w, \\
        0 & \text{otherwise}.
    \end{cases}$$
\end{proposition}
\begin{proof}\leanok Follows from the rigidity degeneration axiom by unwinding the associated graded structure.\end{proof}

\begin{proposition}\label{prereq:prop:syn-einfty-mod-lambda}
    \uses{prereq:ax:rigid,prereq:def:lambda}
    \lean{KIP.Synthetic.einfty_nuX_mod_lambda}\leanok
    For $r \ge 2$, the $E_\infty$-page of the synthetic Adams spectral sequence for $\nu X / \lambda^r$ is
    $$E_\infty^{s,t,w}(\nu X / \lambda^r) \cong \begin{cases}
        Z_{r-t+w}^{s,t}(X) / B_{1+t-w}^{s,t}(X) & \text{if } 0 \le t - w < r, \\
        0 & \text{otherwise}.
    \end{cases}$$
\end{proposition}

\subsubsection{Synthetic Lift}\label{sec:prereq-syn-lift}

\begin{axiom}\label{prereq:ax:synthetic-lift}
    \uses{prereq:ax:nu-cofiber}
    \lean{KIP.Synthetic.synthetic_lift}\leanok
    If a map $f: X \to Y$ has Adams filtration $\AF(f) = k$, then there exists a factorization
    $$\xymatrix{
     & \Sigma^{0,-k}\nu Y \ar[d]^{\lambda^k} \\
     \nu X \ar[r]_{\nu f} \ar@{-->}[ru]^{\Sigma^{0,-k}\tilde{f}} & \nu Y
    }$$
    where $\tilde{f}: \Sigma^{0,k}\nu X \to \nu Y$ is called a \emph{synthetic lift} of $f$.
\end{axiom}

\begin{proposition}\label{prereq:prop:syn-distinguished-triangle}
    \uses{prereq:ax:nu-cofiber,prereq:ax:synthetic-lift}
    \lean{KIP.Synthetic.syn_distinguished_triangle}\leanok
    Suppose that $X \xrightarrow{f} Y \xrightarrow{g} Z \xrightarrow{h} \Sigma X$ is a distinguished triangle of spectra with $\AF(h) > 0$, and consequently a short exact sequence on $\HF$-homology
    $$0 \to H_*X \xrightarrow{H_*f} H_*Y \xrightarrow{H_*g} H_*Z \to 0.$$
    Then there exists a distinguished triangle of synthetic spectra
    $$\nu X \xrightarrow{\nu f} \nu Y \xrightarrow{\nu g} \nu Z \xrightarrow{\Sigma^{0,-1}\hat{h}} \Sigma^{0,-1}\nu \Sigma X = \Sigma^{1,0}\nu X$$
    such that $\nu h = \lambda \hat{h}$.
\end{proposition}

\begin{notation}\label{prereq:not:fhat}
    \uses{prereq:prop:syn-distinguished-triangle}
    \lean{KIP.Synthetic.eHat}\leanok
    For a map $f: X \to Y$ which is part of a distinguished triangle $X \xrightarrow{f} Y \xrightarrow{g} Cf \to \Sigma X$, define
    $$e(f) = \begin{cases}
        0 & \text{if } \AF(f) = 0, \\
        1 & \text{if } \AF(f) > 0.
    \end{cases}$$
    When $\AF(f) = 0$, we also denote $\nu f$ by $\hat{f}$. In both cases, we have $\hat{f}: \Sigma^{0,e(f)}\nu X \to \nu Y$ and $\nu f = \lambda^{e(f)}\hat{f}$. Furthermore, $C\hat{f} \simeq \Sigma^{0,-e(g)}\nu Cf$.
\end{notation}

\begin{proposition}\label{prereq:prop:fhat-triangle}
    \uses{prereq:prop:syn-distinguished-triangle,prereq:not:fhat}
    \lean{KIP.Synthetic.fhat_triangle}\leanok
    Suppose that $X \xrightarrow{f} Y \xrightarrow{g} Z \xrightarrow{h} \Sigma X$ is a distinguished triangle with $e(f) + e(g) + e(h) = 1$. Then there is a distinguished triangle of synthetic spectra
    $$\nu X \xrightarrow{\hat{f}} \Sigma^{0,-e(f)}\nu Y \xrightarrow{\hat{g}} \Sigma^{0,-e(f)-e(g)}\nu Z \xrightarrow{\hat{h}} \Sigma^{0,-1}\nu \Sigma X.$$
\end{proposition}
